% !TEX program = pdflatex
\documentclass[11pt,a4paper]{article}
\usepackage[UTF8]{CJKutf8}
\usepackage[top=0.3in,bottom=0.3in,left=0.3in,right=0.3in]{geometry}
\usepackage{enumitem}
\usepackage{titlesec}
\usepackage{tabularx}
\usepackage{hyperref}
\usepackage{xcolor}

% 颜色定义
\definecolor{linkcolor}{RGB}{0,0,0}
\hypersetup{colorlinks=true,urlcolor=linkcolor}

% 段落间距
\setlength{\parindent}{0pt}
\setlength{\parskip}{0pt}

% section样式
\titleformat{\section}{\normalsize\bfseries}{}{0em}{}[\titlerule]
\titlespacing{\section}{0pt}{4pt}{2pt}

% 列表样式
\setlist[itemize]{leftmargin=1.3em,itemsep=0pt,parsep=0pt,topsep=1pt}

\begin{document}
\begin{CJK}{UTF8}{gbsn}

% ========== 头部信息 ==========
\noindent
\begin{minipage}[t]{0.75\textwidth}
\vspace{4pt}
{\LARGE\bfseries 杨卜慈} \\[4pt]
女\enspace|\enspace 2003.05\enspace|\enspace 手机：15176918735\enspace|\enspace 邮箱：boci\_yang25@tju.edu.cn \\[2pt]
求职意向：\textbf{Agent研发工程师 / AI应用研发工程师}
\end{minipage}%
\begin{minipage}[t]{0.25\textwidth}
\vspace{-4pt}
\raggedleft
% 右侧证件照预留，可直接贴照片，或取消下方注释使用电子照
\framebox[2.2cm]{\rule{0pt}{2.8cm}}
% \includegraphics[width=2.2cm,height=2.8cm]{photo.jpg}
\end{minipage}

\vspace{-4pt}

% ========== 教育背景 ==========
\section{教育背景}
\textbf{天津大学（985、211）}\hfill 2025.09 -- 2028.06 \\
电子信息（大模型/Agent方向）硕士 \hfill 主修课程：高级人工智能（92分）、软件体系结构（89分）

\vspace{2pt}
\textbf{河北师范大学}\hfill 2021.09 -- 2025.06 \\
计算机科学与技术 本科 \hfill 绩点表现：GPA 3.8/4.0，加权均分 90.80，专业排名 \textbf{1/141} 

% ========== 个人亮点与核心荣誉 ==========
\section{个人亮点与核心荣誉}
\begin{itemize}
\item \textbf{算法竞赛}： \textbf{CCPC} 中国大学生程序设计竞赛（女生专场）\textbf{铜奖}、蓝桥杯程序设计竞赛国三、天梯赛国三；
\item \textbf{学业荣誉}：本科专业排名 \textbf{1/141}，荣获校级一等奖学金（2023）及连续三年校级三好学生（2021-2023）；
\item \textbf{校园实践}：担任\textbf{天大心理健康中心咨询助理}，对接协调多名专业咨询师排班，具备良好的沟通协作与应急危机接待能力，熟悉标准 SOP 并赋能 Agent 逻辑设计；
\end{itemize}

\section{自我评价}
具备从问题定义、系统设计到工程落地的完整项目经验；擅长将复杂研究问题拆解为可执行的模块化方案，并在迭代中保持代码质量与可观测性；沟通主动，能与非技术角色高效对齐需求边界。


% ========== 专业技能 ==========
\section{专业技能}
\begin{itemize}
\item \textbf{Agent 系统设计：}理解 Agent、Workflow 与 Agent Loop 的边界，具备确定性状态机管理用户流程与人工确认的实践；使用 CrewAI 完成单/多角色编排与受限上下文协作
\item \textbf{多智能体与来源约束：}使用 Python/Pydantic 设计结构化消息与状态契约；具备公开视图隔离、多 Agent 消息传播和来源依赖建模经验，可区分同源重复、独立佐证与实际暴露证据
\item \textbf{可靠性与可观测性：}实现 JSON 输出校验、确定性回退、超时重试、状态快照和事件账本；支持故障定位与轨迹回放
\item \textbf{大模型工程化：}熟悉 OpenAI 兼容接口、Prompt/Context 设计与模型输出约束；使用 FakeProvider/FakeTransport 配合 pytest 完成离线故障注入和回归测试
\item \textbf{Vibe Coding 快速原型：}熟练使用 Codex、Claude 等 AI 编程工具完成从需求拆解、代码生成、调试重构到可部署原型的全流程交付
\end{itemize}

% ========== 项目经历 ==========
\section{项目（科研）经历}

\textbf{EviCon-Lab：可审计的 LLM 多智能体信息传播与来源依赖评测框架} \hfill 2026.03 -- 至今 \\
\textit{项目主导} \enspace|\enspace Python / Pydantic / OpenAI-compatible API / pytest / JSONL \enspace|\enspace \href{https://github.com/Bobozane/evicon-lab}{GitHub}

面向 LLM 多智能体网络中的来源依赖与信息传播可靠性问题，主导设计并实现可审计实验运行时，将来源图、消息暴露、轮次状态与回放恢复统一纳入执行链路。
\begin{itemize}
\item \textbf{来源依赖建模：}设计来源图与数据契约，将同一上游内容的转发归并至统一来源根；构建 Agent 公开视图，限制决策仅使用实际暴露的消息与证据
\item \textbf{受控 Agent 运行时：}基于 Pydantic 约束判断、采纳、分享与引用等结构化输出；实现 OpenAI 兼容 Provider，统一 JSON Schema、超时、重试、异常分类与请求 fingerprint
\item \textbf{传播回放与恢复：}构建 append-only Exposure Ledger、轮次快照及状态恢复机制，支持核查传播路径、干预时序与可见性权限
\item \textbf{确定性工程验证：}使用 FakeProvider/FakeTransport 完成 4 类场景 $\times$ 4 种条件 $\times$ 48 次确定性运行，传播与干预回放 48/48 通过；自动化测试 \textbf{1484 passed, 1 skipped}
\end{itemize}
\vspace{3pt}
\textbf{职业冲突排练台：用户主导的 AI 价值反思 Agent} \hfill 2026.06 -- 至今 \\
\textit{独立设计与开发} \enspace|\enspace Python / CrewAI / Pydantic / Streamlit / SQLite \enspace|\enspace \href{https://github.com/Bobozane/career-value-stage}{GitHub}

\vspace{1pt}
面向职业选择中的价值取舍与信息不确定性，独立设计并实现多阶段 AI 反思 Agent，将价值角色、A/B 取舍、AI 表达校准和现实信息核实组织为可审计的交互流程。
\begin{itemize}
\item \textbf{Agent 工作流与结构化输出：}基于 CrewAI 编排角色提议、对话回应、条件整合和总结；使用 Pydantic 校验 JSON 输出，并为模型失败提供确定性回退
\item \textbf{用户可控状态机：}实现价值角色确认、A/B 立场接地、可选角色互动、接受/修改/拒绝 AI 草案、暂停与跳过；未经用户确认的内容不会进入最终总结
\item \textbf{研究数据与可靠性：}设计 SQLite 事件账本和会话快照，记录角色修订、AI 校准、模型版本、耗时、错误及前后测；支持 Markdown、CSV 和 JSON 导出，并采用匿名编号与幂等保存
\item \textbf{测试与形成性验证：}使用 FakeProvider/FakeTransport 进行离线故障注入和自动化测试；通过 4 位参与者迭代角色互动适用边界、停止机制和现实信息核实出口
\end{itemize}

\end{CJK}
\end{document}
    
